\documentclass[11pt,a4paper]{article}
\usepackage[top=1in,bottom=1in,left=1.1in,right=1.1in]{geometry}
\usepackage{mathptmx}
\usepackage{hyperref}
\pagestyle{empty}
\begin{document}

\begin{center}
{\Large\bfseries Modeling electrical losses in CsPbI3 solar cells: stacked idealizations and a routine-quality envelope at 18.0\%}
\end{center}

\begin{center}
Hussain Touhid Siddiquee\textsuperscript{*,a} and Md.\ Abdul Malek Fahim\textsuperscript{a}
\end{center}

\begin{center}
\textsuperscript{a}Department of EEE, Leading University, Sylhet-3112, Bangladesh
\end{center}

\begin{center}
\textsuperscript{*}Corresponding author. E-mail: \url{touhidsiddiqueeraj@gmail.com}
\end{center}

\vspace{1em}
\noindent\textbf{Article type:} Regular Paper, for consideration in \textit{Solar Energy} (Elsevier).

\vspace{1em}
\noindent\textbf{Abstract.} Stacked idealisations inflate SCAPS-1D CsPbI3 projections. On one fully worked CsPbI3/CBTS stack, joint optimisation of thickness, doping, defects, gap, and both transport-layer affinities gives a permissive-box ceiling of \textbf{24.95\%}. Sampling a frozen routine-quality prior volume (96-point Latin hypercube over eight factors; ranges and measured anchors in Table~1 of the manuscript) gives a median of \textbf{18.0\%} with a 5--95\% prior-predictive envelope of \textbf{15.6--20.0\%}. This is a sensitivity envelope under stated priors, not a confidence interval on real devices. The 2024--2026 experimental records fall above that envelope. The ceiling record comparison is density-of-states conditional: with capped density of states, literature radiative recombination ($C_r=10^{-10}$~cm$^3$/s) lowers the ceiling to 21.86\%, below the 22.02\% certified record; with effective-mass physical density of states, the same radiative term gives 23.46\% ($N_v=4.1\times10^{18}$~cm$^{-3}$) and 23.04\% ($N_v=8.9\times10^{18}$~cm$^{-3}$), above it. Interface sensitivity is regime dependent, bulk defects dominate at routine defect density, and series resistance up to 5~$\Omega$cm$^2$ closes roughly half of the six-point fill-factor gap against measured devices, with the remainder attributed to nonradiative interface recombination (ideality above unity).

\vspace{1em}
\noindent\textbf{Keywords:} CsPbI3; perovskite solar cell; SCAPS-1D; CBTS; reproducibility; device simulation

\vspace{1em}
\noindent\textbf{Author contributions (CRediT).} Hussain Touhid Siddiquee: conceptualization, methodology, software, formal analysis, investigation, data curation, visualization, validation, and writing---original draft, review and editing. Md.\ Abdul Malek Fahim: software, validation, and writing---original draft.

\vspace{1em}
\noindent\textbf{Data availability.} The simulation archive is based on repository revision \texttt{6e991c1}. Definition files, sweep receipts, raw curves, scripts, verification receipts, and typeset sources are in the public repository at \url{https://github.com/touhidsiddiqueeraj-bit/Modeling-electricModeling-electrical-losses-in-CsPbI3-solar-cells-al-losses-in-CsPbI3-solar-cells-}.

\end{document}
